\chapter{2007年安徽卷（文）}

\section{单选题}

\begin{problem}
若 \(A = \{x \mid x^2 = 1\}\)，\(B = \{x \mid x^2 - 2x - 3 = 0\}\)，则 \(A \cap B =\)（\quad）

\choices
    {\(\{3\}\)}
    {\( \{1\} \)}
    {\(\varnothing\)}
    {\(\{-1\}\)}
\end{problem}

\begin{answer}
D
\end{answer}

\begin{solution}
由 \(x^{2}=1\) 得 \(x=1\) 或 \(x=-1\)，所以 \(A=\{-1,1\}\). 因式分解
\(x^{2}-2x-3=(x-3)(x+1)\)，得 \(B=\{-1,3\}\). 因而
\(A\cap B=\{-1\}\)，对应选项 D.
\end{solution}

\begin{problem}
椭圆 \(x^{2} + 4y^{2} = 1\) 的离心率为（\quad）

\choices
    {\(\frac{\sqrt{3}}{2}\)}
    {\(\frac{3}{4}\)}
    {\(\frac{\sqrt{2}}{2}\)}
    {\(\frac{2}{3}\)}
\end{problem}

\begin{answer}
A
\end{answer}

\begin{solution}
将椭圆方程化为标准形式
\(\dfrac{x^{2}}{1^{2}}+\dfrac{y^{2}}{\left(\frac12\right)^{2}}=1\). 因此半长轴
\(a=1\)，半短轴 \(b=\frac12\)，焦半距
\(c=\sqrt{a^{2}-b^{2}}=\frac{\sqrt3}{2}\). 所以离心率
\(e=\dfrac ca=\dfrac{\sqrt3}{2}\)，对应选项 A.
\end{solution}

\begin{problem}
等差数列 \(\{a_{n}\}\) 的前 \(n\) 项和为 \(S_{n}\). 若 \(a_{2} = 1\)，\(a_{3} = 3\)，\(S_{4} =\)（\quad）

\choices
    {\(12\)}
    {\(10\)}
    {\(8\)}
    {\(6\)}
\end{problem}

\begin{answer}
C
\end{answer}

\begin{solution}
设等差数列的公差为 \(d\). 由 \(a_{3}-a_{2}=3-1\) 得 \(d=2\)，从而
\(a_{1}=a_{2}-d=-1\). 因此
\(S_{4}=a_{1}+a_{2}+a_{3}+a_{4}=-1+1+3+5=8\)，对应选项 C.
\end{solution}

\begin{problem}
下列函数中，反函数是其自身的函数为（\quad）

\choices
    {\(f(x) = x^{2}\)（\(x \in [0, +\infty)\)）}
    {\(f(x) = x^{3}\)（\(x \in (-\infty, +\infty)\)）}
    {\(f(x) = \e^{x}\)（\(x \in (-\infty, +\infty)\)）}
    {\(f(x) = \frac{1}{x}\)（\(x \in (0, +\infty)\)）}
\end{problem}

\begin{answer}
D
\end{answer}

\begin{solution}
逐项求反函数. 选项 A 的反函数为 \(y=\sqrt{x}\)，不等于原函数；选项 B 的反函数为
\(y=\sqrt[3]{x}\)，不等于原函数；选项 C 的反函数为 \(y=\ln x\)，定义域也不同；
选项 D 在 \((0,+\infty)\) 上满足 \(f(f(x))=\dfrac{1}{1/x}=x\)，故其反函数就是自身，
对应选项 D.
\end{solution}

\begin{problem}
若圆 \(x^{2} + y^{2} - 2x - 4y = 0\) 的圆心到直线 \(x - y + a = 0\) 的距离为 \(\frac{\sqrt{2}}{2}\)，则 \(a\) 的值为（\quad）

\choices
    {\(-2\) 或 \(2\)}
    {\(\frac{1}{2}\) 或 \(\frac{3}{2}\)}
    {\(2\) 或 \(0\)}
    {\(-2\) 或 \(0\)}
\end{problem}

\begin{answer}
C
\end{answer}

\begin{solution}
将圆的方程配方得
\((x-1)^{2}+(y-2)^{2}=5\)，所以圆心为 \(O(1,2)\). 圆心到直线
\(x-y+a=0\) 的距离为
\[
\frac{|1-2+a|}{\sqrt{1^{2}+(-1)^{2}}}=\frac{|a-1|}{\sqrt2}
\]
由题意 \(\dfrac{|a-1|}{\sqrt2}=\dfrac{\sqrt2}{2}\)，得 \(|a-1|=1\)，所以
\(a=0\) 或 \(a=2\)，对应选项 C.
\end{solution}

\begin{problem}
设 \(l\)，\(m\)，\(n\) 均为直线，其中 \(m\)，\(n\) 在平面 \(\alpha\) 内，“\(l \perp \alpha\)”是“\(l \perp m\) 且 \(l \perp n\)”的（\quad）

\choices
    {充分不必要条件}
    {必要不充分条件}
    {充分必要条件}
    {既不充分也不必要条件}
\end{problem}

\begin{answer}
A
\end{answer}

\begin{solution}
若 \(l\perp\alpha\)，则 \(l\) 垂直于平面 \(\alpha\) 内的直线，因而
\(l\perp m\) 且 \(l\perp n\)，所以“\(l\perp\alpha\)”是该结论的充分条件.

它不是必要条件. 例如取 \(\alpha:z=0\)，取平面内的直线
\(m:y=0\)、\(n:y=1\)，再取平面内的直线 \(l:x=0\). 直线 \(l\) 分别与
\(m\)，\(n\) 垂直，但 \(l\subset\alpha\)，所以 \(l\) 不垂直于 \(\alpha\). 因此对应选项 A.
\end{solution}

\begin{problem}
图中的图象所表示的函数的解析式为（\quad）

\bitmapfigure[max width=0.25\linewidth]{img/2007/anhui_liberal/q7_fig1.png}

\choices
    {\( y = \frac{3}{2} |x - 1| \)（\(0 \leqslant x \leqslant 2\)）}
    {\( y = \frac{3}{2} - \frac{3}{2} |x - 1| \)（\(0 \leqslant x \leqslant 2\)）}
    {\( y = \frac{3}{2} - |x - 1| \)（\(0 \leqslant x \leqslant 2\)）}
    {\( y = 1 - |x - 1| \)（\(0 \leqslant x \leqslant 2\)）}
\end{problem}

\begin{answer}
B
\end{answer}

\begin{solution}
从图象可读出折线经过 \((0,0)\)、\(\left(1,\frac32\right)\)、\((2,0)\). 当
\(0\leqslant x\leqslant1\) 时，直线斜率为 \(\frac32\)，函数为 \(y=\frac32x\)；当
\(1\leqslant x\leqslant2\) 时，函数为 \(y=\frac32(2-x)\). 将两段合写为
\(y=\frac32-\frac32|x-1|\)，定义域为 \(0\leqslant x\leqslant2\)，对应选项 B.
\end{solution}

\begin{problem}
设 \(a > 1\)，且 \(m = \log_a(a^2 + 1)\)，\(n = \log_a(a - 1)\)，\(p = \log_a(2a)\)，则 \(m\)，\(n\)，\(p\) 的大小关系为（\quad）

\choices
    {\(n > m > p\)}
    {\(m > p > n\)}
    {\(m > n > p\)}
    {\(p > m > n\)}
\end{problem}

\begin{answer}
B
\end{answer}

\begin{solution}
因为 \(a>1\)，对数函数 \(\log_a x\) 在正数范围内单调递增. 又
\(a^{2}+1-2a=(a-1)^{2}>0\)，所以 \(a^{2}+1>2a\)，从而 \(m>p\). 同时
\(2a>a-1>0\)，所以 \(p>n\). 因而 \(m>p>n\)，对应选项 B.
\end{solution}

\begin{problem}
如果点 \(P\) 在平面区域 \(\begin{cases}
2x - y + 2 \geqslant 0,\\
x + y - 2 \leqslant 0,\\
2y - 1 \geqslant 0
\end{cases}\) 上，点 \(Q\) 在曲线 \(x^{2} + (y + 2)^{2} = 1\) 上，那么 \( \left|PQ\right| \) 的最小值为（\quad）

\choices
    {\(\frac{3}{2}\)}
    {\(\frac{4}{\sqrt{5}} - 1\)}
    {\(2\sqrt{2} - 1\)}
    {\(\sqrt{2} - 1\)}
\end{problem}

\begin{answer}
A
\end{answer}

\begin{solution}
曲线 \(x^{2}+(y+2)^{2}=1\) 的圆心为 \(O(0,-2)\)，半径为 \(1\). 对任意满足条件的
\(P(x,y)\)，由 \(y\geqslant\frac12\) 得
\[
|PO|=\sqrt{x^{2}+(y+2)^{2}}\geqslant y+2\geqslant\frac52
\]
由三角不等式，任意圆上点 \(Q\) 满足
\(\left|PQ\right|\geqslant\bigl||PO|-|OQ|\bigr|\geqslant\frac52-1=\frac32\).

点 \(P=\left(0,\frac12\right)\) 满足三个不等式，点 \(Q=(0,-1)\) 在该圆上，且
\(\left|PQ\right|=\frac32\). 所以最小值为 \(\frac32\)，对应选项 A.
\end{solution}

\begin{problem}
把边长为 \(\sqrt{2}\) 的正方形 \(ABCD\) 沿对角线 \(AC\) 折成直二面角，折成直二面角后，在 \(A\)，\(B\)，\(C\)，\(D\) 四点所在的球面上，\(B\) 与 \(D\) 两点之间的球面距离为（\quad）

\choices
    {\(\sqrt{2}\pi\)}
    {\(\pi\)}
    {\(\frac{\pi}{2}\)}
    {\(\frac{\pi}{3}\)}
\end{problem}

\begin{answer}
C
\end{answer}

\begin{solution}
折叠后取坐标系，使折痕 \(AC\) 为 \(x\) 轴，并令两个互相垂直的平面分别为
\(z=0\) 与 \(y=0\). 因为 \(|AC|=\sqrt2\cdot\sqrt2=2\)，可取
\(A=(0,0,0)\)，\(C=(2,0,0)\)，\(B=(1,1,0)\)，\(D=(1,0,1)\).
点 \(O=(1,0,0)\) 到四点的距离均为 \(1\)，故这四点所在球面的球心为 \(O\)，半径为
\(R=1\). 又
\(\overrightarrow{OB}=(0,1,0)\)、\(\overrightarrow{OD}=(0,0,1)\)，所以
\(\overrightarrow{OB}\cdot\overrightarrow{OD}=0\)，中心角 \(\angle BOD=\frac\pi2\).
球面距离等于半径乘中心角，故
\(s=R\angle BOD=\frac\pi2\)，对应选项 C.
\end{solution}

\begin{problem}
定义在 \(\R\) 上的函数 \(f(x)\) 既是奇函数，又是周期函数，\(T\) 是它的一个正周期. 若将方程 \(f(x) = 0\) 在闭区间 \([-T, T]\) 上的根的个数记为 \(n\)，则 \(n\) 可能为（\quad）

\choices
    {\(0\)}
    {\(1\)}
    {\(3\)}
    {\(5\)}
\end{problem}

\begin{answer}
D
\end{answer}

\begin{solution}
因为 \(f\) 是奇函数，\(f(0)=0\). 又 \(T\) 是正周期，所以
\(f(-T)=f(0)=0\)、\(f(T)=f(0)=0\). 另外由周期性和奇性，
\[
f\left(\frac T2\right)=f\left(-\frac T2+T\right)=f\left(-\frac T2\right)
= -f\left(\frac T2\right)
\]
故 \(f\left(\frac T2\right)=0\)，同理 \(f\left(-\frac T2\right)=0\). 因而
\(-T\)，\(-\frac T2\)，\(0\)，\(\frac T2\)，\(T\) 都是闭区间 \([-T,T]\) 内的根，根的个数至少为
\(5\). 例如 \(f(x)=\sin\frac{2\pi x}{T}\) 是奇函数且以 \(T\) 为正周期，在该区间内恰有这五个根，
所以 \(n\) 可能为 \(5\)，对应选项 D.
\end{solution}

\section{填空题}

\begin{problem}
已知 \( (1 - x)^{5} = a_{0} + a_{1}x + a_{2}x^{2} + a_{3}x^{3} + a_{4}x^{4} + a_{5}x^{5} \)，则 \( (a_{0} + a_{2} + a_{4})(a_{1} + a_{3} + a_{5}) \) 的值等于\fillinblank{}.
\end{problem}

\begin{answer}
\(-256\)
\end{answer}

\begin{solution}
令 \(P(x)=(1-x)^{5}=a_{0}+a_{1}x+a_{2}x^{2}+a_{3}x^{3}+a_{4}x^{4}+a_{5}x^{5}\). 取
\(x=1\) 和 \(x=-1\)，得
\(P(1)=a_{0}+a_{1}+a_{2}+a_{3}+a_{4}+a_{5}=0\)，以及
\(P(-1)=a_{0}-a_{1}+a_{2}-a_{3}+a_{4}-a_{5}=32\). 因而
\(a_{0}+a_{2}+a_{4}=\frac{P(1)+P(-1)}2=16\)，\(a_{1}+a_{3}+a_{5}=\frac{P(1)-P(-1)}2=-16\).
所求值为 \(16\times(-16)=-256\).
\end{solution}

\begin{problem}
在四面体 \(O - ABC\) 中，\(\overrightarrow{AB} = \bs{a}\)，\(\overrightarrow{OB} = \bs{b}\)，\(\overrightarrow{OC} = \bs{c}\)，\(D\) 为 \(BC\) 的中点，\(E\) 为 \(AD\) 的中点，则 \(\overrightarrow{OE} =\)\fillinblank{}.（用 \(\bs{a}\)，\(\bs{b}\)，\(\bs{c}\) 表示）
\end{problem}

\begin{answer}
\(-\frac12\bs{a}+\frac34\bs{b}+\frac14\bs{c}\)
\end{answer}

\begin{solution}
以 \(O\) 为原点，用位置向量表示各点. 由 \(\overrightarrow{AB}=\bs{a}\) 和
\(\overrightarrow{OB}=\bs{b}\)，得 \(\overrightarrow{OA}=\bs{b}-\bs{a}\)，而
\(\overrightarrow{OC}=\bs{c}\). 因为 \(D\) 是 \(BC\) 的中点，
\[
\overrightarrow{OD}=\frac{\overrightarrow{OB}+\overrightarrow{OC}}2
=\frac12\bs{b}+\frac12\bs{c}
\]
又因为 \(E\) 是 \(AD\) 的中点，
\[
\overrightarrow{OE}=\frac{\overrightarrow{OA}+\overrightarrow{OD}}2
=\frac12\left(\bs{b}-\bs{a}+\frac12\bs{b}+\frac12\bs{c}\right)
=-\frac12\bs{a}+\frac34\bs{b}+\frac14\bs{c}
\]
\end{solution}

\begin{problem}
在正方体上任意选择两条棱，则这两条棱相互平行的概率为\fillinblank{}.
\end{problem}

\begin{answer}
\(\frac3{11}\)
\end{answer}

\begin{solution}
正方体有 \(12\) 条棱，任意选择两条棱的总数为
\(\mathrm C_{12}^{2}=66\). 正方体的棱按方向分为三组，每组有四条互相平行的棱，
所以相互平行的棱对共有
\(3\mathrm C_{4}^{2}=18\) 对. 因此所求概率为
\[
\frac{18}{66}=\frac3{11}
\]
\end{solution}

\begin{problem}
函数 \( f(x)=3\sin\left(2x-\frac{\pi}{3}\right) \) 的图象为 \(C\)，如下结论中正确的是\fillinblank{}.（写出所有正确结论的编号）

\begin{circlelist}
\item 图象 \(C\) 关于直线 \( x = \frac{11\pi}{12} \) 对称；
\item 图象 \(C\) 关于点 \( \left(\frac{2\pi}{3},0\right) \) 对称；
\item 函数 \( f(x) \) 在区间 \( \left(-\frac{\pi}{12}, \frac{5\pi}{12}\right) \) 内是增函数；
\item 由 \(y = 3\sin 2x\) 的图象向右平移 \(\frac{\pi}{3}\) 个单位长度可以得到图象 \(C\).
\end{circlelist}
\end{problem}

\begin{answer}
\(\circled{1}\circled{2}\circled{3}\)
\end{answer}

\begin{solution}
对结论 1，正弦函数的对称轴满足
\(2x-\frac\pi3=\frac\pi2+k\pi\)，即
\(x=\frac{5\pi}{12}+\frac{k\pi}{2}\). 取 \(k=1\) 得
\(x=\frac{11\pi}{12}\)，所以结论 1 正确.

对结论 2，图象的中心对称点的横坐标满足
\(2x-\frac\pi3=k\pi\)，即 \(x=\frac\pi6+\frac{k\pi}{2}\). 取
\(x=\frac{2\pi}{3}\) 时 \(k=1\)，且对于任意 \(h\)，
\[
f\left(\frac{2\pi}{3}+h\right)+f\left(\frac{2\pi}{3}-h\right)=0
\]
所以结论 2 正确.

对结论 3，当 \(x\in\left(-\frac\pi{12},\frac{5\pi}{12}\right)\) 时，
\(2x-\frac\pi3\in\left(-\frac\pi2,\frac\pi2\right)\)，从而
\(f'(x)=6\cos\left(2x-\frac\pi3\right)>0\)，所以结论 3 正确.

将 \(y=3\sin2x\) 的图象向右平移 \(\frac\pi3\) 后得到
\(y=3\sin\left(2x-\frac{2\pi}{3}\right)\)，而不是
\(y=3\sin\left(2x-\frac\pi3\right)\). 实际上应向右平移 \(\frac\pi6\)，所以结论 4 错误.
故正确结论的编号为 \(1\)，\(2\)，\(3\).
\end{solution}

\section{解答题}

\begin{problem}
解不等式：\( (|3x-1|-1)(\sin x-2)>0 \).
\end{problem}

\begin{solution}
\(\sin x\leqslant 1\)，所以 \(\sin x-2<0\). 因而原不等式等价于 \(\lvert 3x-1\rvert-1<0\)，即 \(\lvert 3x-1\rvert<1\). 所以 \(-1<3x-1<1\)，解得 \(0<x<\frac{2}{3}\). 故原不等式的解集为 \(\left(0,\frac{2}{3}\right)\).
\end{solution}

\begin{problem}
如图，在六面体 \(ABCD - A_{1}B_{1}C_{1}D_{1}\) 中，四边形 \(ABCD\) 是边长为 2 的正方形，四边形 \(A_{1}B_{1}C_{1}D_{1}\) 是边长为 1 的正方形，\(DD_{1} \perp\) 平面 \(A_{1}B_{1}C_{1}D_{1}\)，\(DD_{1} \perp\) 平面 \(ABCD\)，\(DD_{1} = 2\).

\begin{enumerate}
\item 求证：\(A_{1}C_{1}\) 与 \(AC\) 共面，\(B_{1}D_{1}\) 与 \(BD\) 共面；
\item 求证：平面 \(A_{1}ACC_{1} \perp\) 平面 \(B_{1}BDD_{1}\)；
\item 求二面角 \(A - BB_{1} - C\) 的大小.（用反三角函数值表示）
\end{enumerate}

\bitmapfigure[max width=0.28\linewidth]{img/2007/anhui_liberal/q17_fig1.png}
\end{problem}

\begin{solution}
\begin{enumerate}
\item 建立空间直角坐标系，取 \(D\) 为原点，令 \(DC\)，\(DA\)，\(DD_{1}\) 分别为 \(x\)，\(y\)，\(z\) 轴的正方向. 由题意可取
\(D(0,0,0)\)，\(C(2,0,0)\)，\(A(0,2,0)\)，\(B(2,2,0)\)，\(D_{1}(0,0,2)\)，\(C_{1}(1,0,2)\)，\(A_{1}(0,1,2)\)，\(B_{1}(1,1,2)\). 因此 \(\overrightarrow{A_{1}C_{1}}=(1,-1,0)=\frac12\overrightarrow{AC}\)，\(\overrightarrow{B_{1}D_{1}}=(-1,-1,0)=\frac12\overrightarrow{BD}\)，所以 \(A_{1}C_{1}\parallel AC\)，\(B_{1}D_{1}\parallel BD\)，从而 \(A_{1}C_{1}\) 与 \(AC\) 共面，\(B_{1}D_{1}\) 与 \(BD\) 共面.
\item 在所取坐标系中，平面 \(B_{1}BDD_{1}\) 内的两条方向向量可取
\(\overrightarrow{DB}=(2,2,0)\) 和 \(\overrightarrow{DD_{1}}=(0,0,2)\). 向量
\(\bs{n}_{2}=(1,-1,0)\) 同时垂直于它们，因而是平面 \(B_{1}BDD_{1}\) 的一个法向量.
又 \(\overrightarrow{AC}=(2,-2,0)=2\bs{n}_{2}\)，所以 \(AC\perp\) 平面
\(B_{1}BDD_{1}\). 由于 \(AC\subset\) 平面 \(A_{1}ACC_{1}\)，可得平面
\(A_{1}ACC_{1}\perp\) 平面 \(B_{1}BDD_{1}\).
\item 设 \(\bs{u}=\overrightarrow{BB_{1}}=(-1,-1,2)\)，则 \(\bs{u}\cdot\bs{u}=6\). 过 \(B\) 作平面垂直于 \(BB_{1}\)，将 \(\overrightarrow{BA}\) 与 \(\overrightarrow{BC}\) 投影到这个平面上，所得两个方向向量分别为
\[
\bs{p}=\overrightarrow{BA}-\frac{\overrightarrow{BA}\cdot\bs{u}}{\bs{u}\cdot\bs{u}}\bs{u}
=(-2,0,0)-\frac13(-1,-1,2)
=\left(-\frac53,\frac13,-\frac23\right)
\]
和
\[
\bs{q}=\overrightarrow{BC}-\frac{\overrightarrow{BC}\cdot\bs{u}}{\bs{u}\cdot\bs{u}}\bs{u}
=(0,-2,0)-\frac13(-1,-1,2)
=\left(\frac13,-\frac53,-\frac23\right)
\]
这两个方向所成的角就是二面角 \(A-BB_{1}-C\) 的平面角 \(\theta\). 于是 \(\bs{p}\cdot\bs{q}=-\frac23\)，且 \(|\bs{p}|^{2}=|\bs{q}|^{2}=\frac{10}{3}\)，所以
\[
\cos\theta=\frac{\bs{p}\cdot\bs{q}}{|\bs{p}||\bs{q}|}
=\frac{-\frac23}{\frac{10}{3}}=-\frac15
\]
故二面角 \(A-BB_{1}-C\) 的大小为 \(\arccos\left(-\frac15\right)\).
\end{enumerate}
\end{solution}

\begin{problem}
设 \(F\) 是抛物线 \(G: x^{2} = 4y\) 的焦点.

\begin{enumerate}
\item 过点 \( P(0, -4) \) 作抛物线 \(G\) 的切线，求切线方程；
\item 设 \(A\)，\(B\) 为抛物线 \(G\) 上异于原点的两点，且满足 \( \overrightarrow{FA} \cdot \overrightarrow{FB} = 0 \)，延长 \(AF\)，\(BF\) 分别交抛物线 \(G\) 于点 \(C\)，\(D\)，求四边形 \(ABCD\) 面积的最小值.
\end{enumerate}
\end{problem}

\begin{solution}
\begin{enumerate}
\item 设切点为 \(X(2t,t^{2})\). 由 \(x^{2}=4y\) 得 \(y'=x/2\)，所以在 \(X\) 处切线的斜率为 \(t\)，切线方程为 \(y-t^{2}=t(x-2t)\)，即 \(y=tx-t^{2}\). 将 \(P(0,-4)\) 代入，得 \(-4=-t^{2}\)，所以 \(t=\pm2\). 因而所求切线方程为 \(y=2x-4\) 或 \(y=-2x-4\).
\item 抛物线的焦点为 \(F(0,1)\). 设 \(A(2u,u^{2})\)，\(B(2v,v^{2})\)，其中 \(u\neq0\)，\(v\neq0\). 对于参数为 \(s\) 的点 \(X(2s,s^{2})\)，点 \(F\)，\(A\)，\(X\) 共线的条件为
\[
u(s^{2}-1)-s(u^{2}-1)=(s-u)(us+1)=0
\]
除去 \(s=u\) 对应的点 \(A\)，直线 \(AF\) 与抛物线的另一交点 \(C\) 的参数为 \(-1/u\). 同理，\(D\) 的参数为 \(-1/v\).

由于
\(|FA|=\sqrt{4u^{2}+(u^{2}-1)^{2}}=u^{2}+1\)，\(|FC|=\sqrt{\frac4{u^{2}}+\left(\frac1{u^{2}}-1\right)^{2}}=\frac{u^{2}+1}{u^{2}}\)
且 \(A\)，\(F\)，\(C\) 三点依次在同一直线上，所以
\[
|AC|=\frac{(u^{2}+1)^{2}}{u^{2}}
\]
同理 \(|BD|=(v^{2}+1)^{2}/v^{2}\). 由 \(\overrightarrow{FA}\cdot\overrightarrow{FB}=0\)，直线 \(AC\) 与 \(BD\) 垂直，因此四边形 \(ABCD\) 的面积为
\[
S=\frac12|AC||BD|
\]
令
\(p_{1}=\frac{2u}{u^{2}+1}\)，\(q_{1}=\frac{u^{2}-1}{u^{2}+1}\)，\(p_{2}=\frac{2v}{v^{2}+1}\)，\(q_{2}=\frac{v^{2}-1}{v^{2}+1}\).
则 \(p_{1}^{2}+q_{1}^{2}=p_{2}^{2}+q_{2}^{2}=1\)，而已知条件给出 \(p_{1}p_{2}+q_{1}q_{2}=0\). 两个单位向量垂直，所以 \(p_{2}^{2}=q_{1}^{2}\)，从而 \(p_{1}^{2}+p_{2}^{2}=1\). 又
\(|AC|=\frac4{p_{1}^{2}}\)，\(|BD|=\frac4{p_{2}^{2}}\)
故
\[
S=\frac8{p_{1}^{2}p_{2}^{2}}
\geqslant\frac8{\left(\frac{p_{1}^{2}+p_{2}^{2}}2\right)^{2}}
=32
\]
当 \(p_{1}^{2}=p_{2}^{2}=\frac12\) 时等号成立，例如取 \(u=\sqrt2-1\)，\(v=\sqrt2+1\). 所以四边形 \(ABCD\) 面积的最小值为 \(32\).
\end{enumerate}
\end{solution}

\begin{problem}
在医学生物学试验中，经常以果蝇作为试验对象，一个关有 \(6\) 只果蝇的笼子里，不慎混入了两只苍蝇（此时笼内共有 \(8\) 只蝇子：\(6\) 只果蝇和 \(2\) 只苍蝇），只好把笼子打开一个小孔，让蝇子一只一只地往外飞，直到两只苍蝇都飞出，再关闭小孔.

\begin{enumerate}
\item 求笼内恰好剩下 \(1\) 只果蝇的概率；
\item 求笼内至少剩下 \(5\) 只果蝇的概率.
\end{enumerate}
\end{problem}

\begin{solution}
\begin{enumerate}
\item 只考虑两只苍蝇在依次飞出顺序中的位置. 这两个位置从 \(8\) 个位置中任取，所有可能情形共有 \(\mathrm C_{8}^{2}\) 种，且等可能. 设两只苍蝇的位置为 \(i<j\)，则第 \(j\) 只飞出的是第二只苍蝇，此时笼内剩下的果蝇数为 \(8-j\). 恰好剩下 \(1\) 只果蝇等价于 \(j=7\)，有 \(6\) 种选择 \(i\)，所以所求概率为 \(\frac{6}{\mathrm C_{8}^{2}}=\frac3{14}\).
\item 笼内至少剩下 \(5\) 只果蝇等价于 \(8-j\geqslant5\)，即 \(j\leqslant3\). 这等价于两只苍蝇的位置都在前 \(3\) 个位置中，共有 \(\mathrm C_{3}^{2}\) 种情形，所以所求概率为 \(\frac{\mathrm C_{3}^{2}}{\mathrm C_{8}^{2}}=\frac3{28}\).
\end{enumerate}
\end{solution}

\begin{problem}
设函数 \(f(x) = -\cos^2 x - 4t\sin \frac{x}{2}\cos \frac{x}{2} + 4t^3 + t^2 - 3t + 4\)，\(x \in \R\)，其中 \( |t| \leqslant 1 \). 将 \( f(x) \) 的最小值记为 \( g(t) \).

\begin{enumerate}
\item 求 \( g(t) \) 的表达式；
\item 讨论 \(g(t)\) 在区间 \((-1, 1)\) 内的单调性并求极值.
\end{enumerate}
\end{problem}

\begin{solution}
\begin{enumerate}
\item 利用 \(\cos^{2}x=1-\sin^{2}x\) 以及 \(2\sin\frac{x}{2}\cos\frac{x}{2}=\sin x\)，有
\[
f(x)=\sin^{2}x-2t\sin x+4t^{3}+t^{2}-3t+3
=(\sin x-t)^{2}+4t^{3}-3t+3
\]
当 \(|t|\leqslant1\) 时，可以取到 \(\sin x=t\)，故 \(f(x)\) 的最小值为
\(g(t)=4t^{3}-3t+3\)，\(-1\leqslant t\leqslant1\).
\item \(g'(t)=12t^{2}-3=3(2t-1)(2t+1)\). 因此 \(g'(t)>0\) 在 \((-1,-\frac12)\) 和 \((\frac12,1)\) 内成立，\(g'(t)<0\) 在 \((-\frac12,\frac12)\) 内成立. 所以 \(g(t)\) 在 \((-1,-\frac12)\) 和 \((\frac12,1)\) 内单调递增，在 \((-\frac12,\frac12)\) 内单调递减.

又 \(g(-\frac12)=4\)，\(g(\frac12)=2\)，所以 \(g(t)\) 在 \((-1,1)\) 内的极大值为 \(4\)，在 \(t=-\frac12\) 处取得；极小值为 \(2\)，在 \(t=\frac12\) 处取得.
\end{enumerate}
\end{solution}

\begin{problem}
某国采用养老储备金制度. 公民在就业的第一年就交纳养老储备金，数目为 \(a_1\)；以后每年交纳的数目均比上一年增加 \(d\)（\(d > 0\)）. 因此，历年所交纳的储备金数目 \(a_1\)，\(a_2\)，\(\cdots\) 是一个公差为 \(d\) 的等差数列. 与此同时，国家给予优惠的计息政策，不仅采用固定利率，而且计算复利. 这就是说，如果固定年利率为 \(r\)（\(r > 0\)），那么，在第 \(n\) 年末，第一年所交纳的储备金就变为 \(a_1(1 + r)^{n-1}\)；第二年所交纳的储备金就变为 \(a_2(1 + r)^{n-2}\)，\(\cdots\). 以 \(T_n\) 表示到第 \(n\) 年末所累计的储备金总额.

\begin{enumerate}
\item 写出 \( T_{n} \) 与 \( T_{n-1} \) （ \( n \geqslant 2 \) ）的递推关系式；
\item 求证：\(T_{n} = A_{n} + B_{n}\). 其中 \(\{A_{n}\}\) 是一个等比数列，\(\{B_{n}\}\) 是一个等差数列.
\end{enumerate}
\end{problem}

\begin{solution}
\begin{enumerate}
\item 令 \(q=1+r\)，则第 \(n\) 年末以前各年缴纳的储备金经过计息后的总额为
\(T_{n}=a_{1}q^{n-1}+a_{2}q^{n-2}+\cdots+a_{n}\)，而
\(qT_{n-1}=a_{1}q^{n-1}+a_{2}q^{n-2}+\cdots+a_{n-1}q\). 因为 \(a_n=a_1+(n-1)d\)，所以
\(T_n=qT_{n-1}+a_n=(1+r)T_{n-1}+a_1+(n-1)d\)，\(n\geqslant2\).
\item 取
\(C=\frac{a_1r+d}{r^{2}}\)，\(A_n=C(1+r)^n\)，\(B_n=-\frac{d}{r}n-C\).
由于 \(1+r\) 为常数，\(\{A_n\}\) 是公比为 \(1+r\) 的等比数列；又 \(B_{n+1}-B_n=-d/r\) 为常数，\(\{B_n\}\) 是等差数列.

下面证明 \(T_n=A_n+B_n\). 当 \(n=1\) 时，
\[
A_1+B_1=C(1+r)-\frac{d}{r}-C=Cr-\frac{d}{r}=a_1=T_1
\]
若 \(T_{n-1}=A_{n-1}+B_{n-1}\)，则由第（1）问的递推关系式，
\[
T_n=(1+r)(A_{n-1}+B_{n-1})+a_1+(n-1)d
=A_n+(1+r)\left(-\frac{d}{r}(n-1)-C\right)+a_1+(n-1)d
\]
由 \(Cr=a_1+\frac{d}{r}\) 得
\(a_1-(1+r)C=-C-\frac{d}{r}\)，又
\(-(1+r)\frac{d}{r}(n-1)+(n-1)d=-\frac{d}{r}(n-1)\)，所以
\[
T_n=A_n-\frac{d}{r}(n-1)-C-\frac{d}{r}=A_n-\frac{d}{r}n-C=A_n+B_n
\]
由数学归纳法，\(T_n=A_n+B_n\) 对所有正整数 \(n\) 成立.
\end{enumerate}
\end{solution}
